../cictikz/src/cictikz/data/tex/cictikz_preamble.tex

% Anti-alias filtering: the same picture as l5_sh with a low pass in front of
% the sample and hold, so the tones outside the Nyquist band never reach it.
%
% Second of three. Geometry comes from tikz/spec_lib.tex, so the axes, the
% wanted signal and the tone positions are identical to l5_sh and l5_subsample
% and the three slides can be compared directly.
%
% The after panel carries no tones at all. That is the point: the filter
% removed them before sampling, so there is nothing left to fold in. Contrast
% with l5_sh, where the same four tones land inside the band.

\begin{circuitikz}[american, transform shape]

  ../cictikz/src/cictikz/data/tex/spec_lib.tex

  % ---- Before sampling: the low pass passes the signal, not the tones ----
  \specAxis{\specYtop}
  % Flat top matched to the signal band, so the magenta tones at 0.6 fs sit
  % visibly outside the skirt rather than on it.
  \specBand{0}{1.05}{1.70}{\specYtop}
  \specBump{\specYtop}
  \specTone{toneA}{-2.64}{\specYtop}     % -1.1 fs
  \specTone{toneB}{-1.44}{\specYtop}     % -0.6 fs
  \specTone{toneB}{1.44}{\specYtop}
  \specTone{toneA}{2.64}{\specYtop}
  \node[specfilter, anchor=east] at (-1.55,\specYtop+1.95) {LP filter};
  \draw[specfilter, thick, ->] (-1.5,\specYtop+1.90) -- (-1.15,\specYtop+1.75);
  \node[anchor=west] at (1.45,\specYtop+1.62) {Before sampling};

  % ---- After sampling: the tones are gone, so nothing folds in ----
  \specAxis{\specYbot}
  \specBump{\specYbot}
  \node[anchor=west] at (1.45,\specYbot+1.62) {After sampling};

  % ---- Low pass, then sample and hold ----
  \specWire{4.2}{4.9}{2.0}
  \specBlock{4.9}{2.0}{LP}
  \specWire{6.4}{7.0}{2.0}
  \specBlock{7.0}{2.0}{S/H}
  \specWire{8.5}{9.2}{2.0}

  \node[anchor=south] at (5.65,3.3) {Anti-alias};
  \draw[thick, ->] (5.65,3.25) -- (5.65,2.55);

\end{circuitikz}
\end{document}
